% Generated by scripts/assemble.py; corrections belong in source/corrections.json.
\examhead{194}{云南大学 2025 年数学分析试卷}{183}
% p184-b18; source PDF page 184
\par\Needspace{5\baselineskip}
\blockanchor{p184-b18}
\majorquestion{一}填空题.\par
% p184-b19; source PDF page 184
\blockanchor{p184-b19}
\question{1}极限  $\lim_{x\to3}\frac{\sqrt{x^{3}+9}-6}{2-\sqrt{x^{3}-23}}=$  \_\_\_\_.\par
% p184-b20; source PDF page 184
\blockanchor{p184-b20}
\question{2}积分 $\int_{2}^{+\infty} \frac{5}{x^4 + 3x^3 - 4} \, \mathrm{d}x =$ \_\_\_\_.\par
% p184-b21; source PDF page 184
\blockanchor{p184-b21}
\question{3}曲面 $z = \frac{x^2}{2} + y^2 - 2$ 的平行于平面 $2x + 2y - z = 0$ 的切平面方程是 \_\_\_\_.\par
% p184-b22; source PDF page 184
\blockanchor{p184-b22}
\question{4}设连续函数 $f(x)$ 满足 $f(x + 2) - f(x) = x$, $\int_0^2 f(x) \, \mathrm{d}x = 0$, 则 $\int_1^3 f(x) \, \mathrm{d}x =$ \_\_\_\_.\par
% p184-b23; source PDF page 184
\blockanchor{p184-b23}
\question{5}设 $F(x) = \int_{0}^{x} \left( \int_{0}^{y^2} \frac{\sin t}{1 + t^2} \, \mathrm{d}t \right) \, \mathrm{d}y$, 则 $F''(x) =$ \_\_\_\_.\par
% p185-b00; source PDF page 185
\blockanchor{p185-b00}
\question{6}级数  $\sum_{n=0}^{\infty}\frac{x^{2n}}{(2n)!}$  的和函数是 \_\_\_\_.\par
% p185-b01; source PDF page 185
\par\Needspace{5\baselineskip}
\blockanchor{p185-b01}
\majorquestion{二}设数列  $\{x_{n}\}$  满足:  $x_{1}>0, x_{n}e^{x_{n+1}}=e^{x_{n}}-1, n=1,2,\cdots$\par
% p185-b02; source PDF page 185
\blockanchor{p185-b02}
(1) 证明数列 $\{x_{n}\}$ 收敛.\par
% p185-b03; source PDF page 185
\blockanchor{p185-b03}
(2) 求 $\lim_{n\to \infty}x_n$\par
% p185-b04; source PDF page 185
\par\Needspace{5\baselineskip}
\blockanchor{p185-b04}
\majorquestion{三}若函数 $f(x)$ 在 $[a, b]$ 上连续且单调增加. 证明: $\int_{a}^{b} x f(x) \mathrm{d}x \geq \frac{a + b}{2} \int_{a}^{b} f(x) \mathrm{d}x$.
\majorquestion{四}已知 $f(x) = \int_{1}^{x} e^{u^2} \mathrm{d}u$.\par
% p185-b05; source PDF page 185
\blockanchor{p185-b05}
\question{1}证明: 存在 $\xi \in (1,2)$, 使得 $f(\xi) = (2 - \xi)e^{\xi^2}$.\par
% p185-b06; source PDF page 185
\blockanchor{p185-b06}
\question{2}证明: 存在 $\eta \in (1,2)$, 使得 $f(2) = \ln 2 \cdot \eta \cdot e^{\eta^2}$.\par
% p185-b07; source PDF page 185
\par\Needspace{5\baselineskip}
\blockanchor{p185-b07}
\majorquestion{五}讨论当 $p$ 取得何值时, 级数 $\sum_{n=1}^{\infty} \frac{(-1)^n}{n^{p+\frac{1}{n}}}$ 发散? 条件收敛? 绝对收敛?
\majorquestion{六}设 $f_n(x) = x n^\alpha e^{-nx}, n = 1,2,\cdots$.\par
% p185-b08; source PDF page 185
\blockanchor{p185-b08}
\question{1}当 $\alpha$ 为何值时, 函数序列 $\{f_n(x)\}$ 在 $[0,1]$ 上收敛?\par
% p185-b09; source PDF page 185
\blockanchor{p185-b09}
\question{2}当 $\alpha$ 为何值时, 函数序列 $\{f_n(x)\}$ 在 $[0,1]$ 上一致收敛?\par
% p185-b10; source PDF page 185
\blockanchor{p185-b10}
\question{3}当 $\alpha$ 为何值时, $\lim_{n\to \infty}\int_0^1 f_n(x)\mathrm{d}x$ 和 $\int_0^1\lim_{n\to \infty}f_n(x)\mathrm{d}x$ 不相等?\par
% p185-b11; source PDF page 185
\par\Needspace{5\baselineskip}
\par\Needspace{6\baselineskip}
\blockanchor{p185-b11}
\majorquestion{七}设\par
\nopagebreak[4]
% p185-b12; source PDF page 185
\blockanchor{p185-b12}
\displayformula{f (x, y) = \left\{ \begin{array}{l l} x y \sin \frac {1}{x ^ {2} + y ^ {2}}, & x ^ {2} + y ^ {2} \neq 0; \\ 0, & x ^ {2} + y ^ {2} = 0. \end{array} \right.}
% p185-b13; source PDF page 185
\blockanchor{p185-b13}
\question{1}证明 $f(x, y)$ 的两个一阶偏导数在点 (0,0) 不连续.\par
% p185-b14; source PDF page 185
\blockanchor{p185-b14}
\question{2}证明 $f(x, y)$ 在 $(0,0)$ 可微.\par
% p185-b15; source PDF page 185
\par\Needspace{5\baselineskip}
\par\Needspace{6\baselineskip}
\blockanchor{p185-b15}
\majorquestion{八}计算三重积分\par
\nopagebreak[4]
% p185-b16; source PDF page 185
\blockanchor{p185-b16}
\displayformula{I = \iiint_ {\Omega} z \mathrm{d} x \mathrm{d} y \mathrm{d} z.}
% p185-b17; source PDF page 185
\par\Needspace{6\baselineskip}
\blockanchor{p185-b17}
其中积分区域 $\Omega$ 是球面 $x^{2} + y^{2} + z^{2} = 8$ 的上半部分和抛物面 $z = \frac{1}{2} (x^2 +y^2)$ 所围立体区域.
\majorquestion{九}设 $L$ 是曲面 $\Sigma :4x^{2} + y^{2} + z^{2} = 1,x\geq 0,y\geq 0,z\geq 0$ 的边界，曲面方向朝上，曲线 $L$ 的方向和曲面的方向符合右手法则，求曲线积分\par
\nopagebreak[4]
% p185-b18; source PDF page 185
\blockanchor{p185-b18}
\displayformula{I = \oint_ {L} (y z ^ {2} - \cos z) \mathrm{d} x + 2 x z ^ {2} \mathrm{d} y + (2 x y z + x \sin z) \mathrm{d} z.}
